\kafkasection{
If cooling behind a shock front is very effective, the gas temperature returns quickly to its preshock value. If we are not interested in the details of the cooling region, we can regard it and the shock itself as a single discontinuity, usually called an isothermal shock, and can treat it in the manner used for adiabatic shocks. Thus, we can retain the mass and momentum jump conditions but replace the energy condition by $T_1 = T_2$, where $T_1$ and $T_2$ are the gas temperatures on each side of the discontinuity, or equivalently the equality of sound speeds: $c_1=c_2=c$. Where $c_1 = (P_1/\rho_1)^{1/2}$ and so forth.
Show that: 
$$v_2 = \frac{c}{\mathcal{M}_\infty}$$
$$\rho_2 = \mathcal{M}^2_\infty \rho_1$$
and,
$$P_2 = \rho_1v^2_1$$
where $\mathcal{M}_\infty$, is the preshock Mach number. Strong isothermal shocks result in very large compressions, and in all cases the postshock gas pressure is equal to the preshock ram pressure. (From Frank, King, \& Raine, Accretion Power in Astrophysics, 2002.)
}
Let's use the sound speed relation with the pressure and density.
\begin{equation}
    P = \rho c_s^2
\end{equation}
Let's also use conservation of mass.
\begin{equation}
    \rho_1 v_1 = \rho_2 v_2
\end{equation}
\begin{equation}
    v_2 = \frac{v_1}{(\rho_2/\rho_1)}
\end{equation}
Finally let's start with the momentum flux.
\begin{equation}
    P_1 + \rho_1 v_1^2 = P_2 + \rho_2 v_2^2
\end{equation}
Using \(P=\rho c_s^2\),
\begin{equation}
    \rho_1 c_s^2 + \rho_1 v_1^2 = \rho_2 c_s^2 + \rho_2 v_2^2 
\end{equation}
Divide through by \(\rho_1\) and substitute \(\rho_2 = (\rho_2/\rho_1)\rho_1\) and
\(v_2 = v_1/(\rho_2/\rho_1)\):
\begin{equation}
c_s^2 + v_1^2
=
(\rho_2/\rho_1) c_s^2 + (\rho_2/\rho_1)\left(\frac{v_1}{(\rho_2/\rho_1)}\right)^2
=
(\rho_2/\rho_1) c_s^2 + \frac{v_1^2}{(\rho_2/\rho_1)}
\end{equation}

Multiply through by \((\rho_2/\rho_1)\):
\begin{equation}
    (\rho_2/\rho_1)(c_s^2 + v_1^2) = (\rho_2/\rho_1)^2 c_s^2 + v_1^2
\end{equation}
Rearranging,
\begin{equation}
    0 = ((\rho_2/\rho_1)^2 - (\rho_2/\rho_1))c_s^2 + (1 - (\rho_2/\rho_1))v_1^2
= ((\rho_2/\rho_1) - 1)((\rho_2/\rho_1) c_s^2 - v_1^2)
\end{equation}

For a non-trivial shock \((\rho_2/\rho_1) \neq 1\), so
\begin{equation}
    (\rho_2/\rho_1) c_s^2 = v_1^2,
\qquad
(\rho_2/\rho_1) = \frac{v_1^2}{c_s^2}
\end{equation}

Introducing the preshock Mach number
\begin{equation}
    \mathcal M_\infty \equiv \frac{v_1}{c_s},
\end{equation}

we get the second relation.
\begin{equation}
    (\rho_2/\rho_1) = \mathcal M_\infty^2.
\end{equation}

Using \(v_2 = v_1/(\rho_2/\rho_1)\),
\begin{equation}
    v_2 = \frac{v_1}{\mathcal M_\infty^2}
= \frac{\mathcal M_\infty c_s}{\mathcal M_\infty^2}
= \frac{c_s}{\mathcal M_\infty}
\end{equation}
which is the first relation. So now we have the first two relations.
\[
\boxed{v_2 = \frac{c_s}{\mathcal M_\infty}}, \qquad
\boxed{\rho_2 = \mathcal M_\infty^2 \rho_1}
\]
Let's use the sound speed and our previous results for the last relation.
\begin{equation}
    P_2 = \rho_2 c_s^2
\end{equation}
\begin{equation}
    P_2 = (\mathcal{M}_\infty^2 \rho_1)c_s^2 = \rho_1 (\mathcal{M}_\infty^2 c_s^2)
\end{equation}
Using the definition of the preshock mach speed:
\begin{equation*}
    \boxed{P_2 = \rho_1 v_1^2}
\end{equation*}
d